\documentclass{standalone}
\usepackage[siunitx, american]{circuitikz}
\usepackage{tikz}
\begin{document}
    \begin{tikzpicture}[>=latex]
        \draw[style=help lines] (0,0) (3,2);
        %Intensidad
        \coordinate (vec1) at (-14:1); 
        %Tensión
        \coordinate (vec2) at (0:1.46);
        %Axis 
        \coordinate (vec3) at (0:3);
        \coordinate (vec4) at (90:3);
        \coordinate (vec5) at (270:3);
        \coordinate (vec6) at (180:3);  

        \draw[->,thick,black] (0,0) -- (vec3) node [below] {$Re$};
        \draw[->,thick,black] (0,0) -- (vec4) node [left] {$Im$};
        \draw[->,thick,black] (0,0) -- (vec5);
        \draw[->,thick,black] (0,0) -- (vec6);  

        \draw[->,thick,red] (0,0) -- (vec1) node[below right] {$\mathrm{\vec{I}}=\SI{0.36}{\ampere}\angle \SI{-14.01}{\degree}$};
        \draw[->,thick,blue] (0,0) -- (vec2) node[above right] {$\mathrm{\vec{V}}=\SI{99.06}{\volt}\angle \SI{0}{\degree}$};

        \draw [red, thick ] (0.9,0) arc (0:-14:1cm)    node [midway, right] {$\varphi$};         
    \end{tikzpicture}
\end{document}